\documentclass[conference]{IEEEtran}
\IEEEoverridecommandlockouts

\usepackage{cite}
\usepackage{amsmath,amssymb,amsfonts}
\usepackage{graphicx}
\usepackage{booktabs}
\usepackage{xcolor}
\usepackage{hyperref}
\usepackage{url}
\usepackage{microtype}
\usepackage{enumitem}

\hypersetup{colorlinks=true,urlcolor=blue,linkcolor=black,citecolor=black}

\begin{document}

\title{Surrogate Aerodynamic Modeling:\\
Predicting Transonic Airfoil Coefficients via Machine Learning}

\author{\IEEEauthorblockN{Ashish Dasu}
\IEEEauthorblockA{Khoury College of Computer Sciences\\
Northeastern University\\
Boston, MA, USA\\
\texttt{dasu.a@northeastern.edu}}}

\maketitle

\begin{abstract}
% TODO (Phase 9). Cover: CFD is slow; transonic regime is hard because of
% shock-induced nonlinearity in Cm; dataset (1362 CFD runs, 8 airfoils,
% Cl+Cm); 4-model comparison; headline result; RAE2822 generalization;
% ~N-x latency speedup over the ~7 min CFD baseline.
Placeholder.
\end{abstract}

\begin{IEEEkeywords}
Machine learning, regression, aerodynamics, transonic flow, surrogate
modeling, airfoil, computational fluid dynamics.
\end{IEEEkeywords}

% ────────────────────────────────────────────────────────────────────
\section{Introduction}
% TODO Phase 9. Motivation (design iteration and GNC demand thousands of
% coefficient evaluations); why transonic is the hard regime; scope of
% this work (Cl, Cm, 8 airfoils, 4 models); contributions vs the source
% paper \cite{elrefaie2024}: polynomial baseline, RAE2822 held-out
% generalization probe, Mach-binned residuals, physical-plausibility
% overlays, learning curves, latency benchmark.

% ────────────────────────────────────────────────────────────────────
\section{Background and Related Work}
% TODO Phase 9. Definitions of Cl, Cm. Why Cm is the hard target in
% transonic flow (shock movement -> aerodynamic center shift). Related
% ML surrogate work: Elrefaie \cite{elrefaie2024}, Bouhlel \cite{bouhlel2020}.

% ────────────────────────────────────────────────────────────────────
\section{Dataset}
% TODO Phase 9. TransonicSurrogate; 1362 rows x 8 airfoils; 18 features
% (16 shape y-coords + alpha + Mach); targets Cl, Cm; airfoil-ID recovery
% by coordinate fingerprinting; 60/20/20 split with RAE2822 held out.

% ────────────────────────────────────────────────────────────────────
\section{Problem Formulation}
% TODO Phase 9. Multi-output regression, joint MSE.

% ────────────────────────────────────────────────────────────────────
\section{Methodology}
\subsection{Polynomial Regression Baseline}
\subsection{Random Forest}
\subsection{XGBoost}
\subsection{Deep Neural Network}
\subsection{Evaluation Protocol}
% 5-fold CV on train; held-out test; RAE2822 unseen-geometry probe;
% Mach-binned residuals; physical plausibility; latency.

% ────────────────────────────────────────────────────────────────────
\section{Results}
\subsection{Test-Set Performance}
\subsection{Generalization to Held-out RAE2822}
\subsection{Residual Analysis by Mach Regime}
\subsection{Physical Plausibility: $C_m$ vs $\alpha$}
\subsection{Feature Importance}
\subsection{Inference Latency}

% ────────────────────────────────────────────────────────────────────
\section{Discussion}
% What worked, what didn't, physical interpretation of failure modes.

% ────────────────────────────────────────────────────────────────────
\section{Limitations}
% Inviscid CFD source data; only 8 airfoils; 2D only; small N.

% ────────────────────────────────────────────────────────────────────
\section{Conclusion}

% ────────────────────────────────────────────────────────────────────
\bibliographystyle{IEEEtran}
\bibliography{refs}

\end{document}
